\section{Phương pháp học sâu}
Các kiến trúc trong phần này được mô tả theo cài đặt tại
\path{deep_models/models.py}. Kích thước tensor được ghi theo thứ tự
$C\times H\times W$ cho CNN; $T\times D$ cho ViT; bỏ qua chiều batch $B$.
Đây là các biến thể cho CIFAR-10, với trọng số được học từ đầu.

\subsection{Nguyên lý chung: học biểu diễn và tối ưu đầu-cuối}
\textbf{Tích chập.} Một lớp convolution tạo mỗi kênh đầu ra bằng cách kết hợp
các vùng lân cận trên tất cả kênh đầu vào:
\begin{equation}
 a_{o,h,w}=b_o+\sum_{c=1}^{C_{in}}\sum_{u,v}
 W_{o,c,u,v}\,x_{c,sh+u-p,sw+v-p}.
\end{equation}
Kernel $W$ được dùng lại ở nhiều vị trí; $s$ là stride, $p$ là padding.
Với kích thước đầu vào $H$, kernel $k$ và dilation 1, kích thước đầu ra là
$H_{out}=\lfloor(H+2p-k)/s\rfloor+1$.
Ví dụ, kernel $3\times3$, padding 1 và stride 1 giữ nguyên kích thước;
stride 2 giảm khoảng một nửa. Việc chia sẻ trọng số giúp cùng một mẫu cạnh
được phát hiện ở nhiều vị trí mà không học một bộ trọng số riêng cho từng pixel.

\textbf{Phi tuyến và gộp thông tin.} ReLU, $\max(0,a)$, làm mạng có khả năng
biểu diễn quan hệ phi tuyến. Nếu chỉ ghép các lớp tuyến tính, toàn mạng vẫn là
một phép biến đổi tuyến tính, dù có nhiều lớp. Max-pooling $2\times2$, stride 2
chọn đáp ứng lớn nhất trong mỗi vùng, giảm kích thước và độ nhạy với dịch chuyển
nhỏ; nó cũng có thể làm mất chi tiết trên ảnh nhỏ. Global average pooling lấy
trung bình trên toàn bộ vùng không gian của từng kênh. Qua nhiều lớp, vùng ảnh
ảnh hưởng đến một đặc trưng (receptive field) mở rộng; biểu diễn có thể kết hợp
mẫu cục bộ thành cấu trúc lớn hơn phục vụ phân lớp.

\textbf{Học biểu diễn.} Mạng phân loại được viết thành
$z=g_\theta(f_\phi(I))$, trong đó $f_\phi$ là bộ trích xuất đặc trưng học được,
$g_\theta$ là bộ phân lớp và $z$ có 10 logit. Với đặc trưng thủ công $h(I)$,
chỉ bộ phân lớp được học; với học sâu, cả $\phi$ và $\theta$ nhận gradient từ
cùng mục tiêu. Chẳng hạn, nếu một mẫu màu hoặc đường biên giúp phân biệt
\emph{ship} với \emph{airplane}, trọng số phía trước có thể được điều chỉnh để
làm nổi bật mẫu đó. Đây là khả năng của mô hình, không phải khẳng định rằng
mọi bộ lọc trong lần chạy này đã được giải thích về mặt ngữ nghĩa.

\textbf{Hàm mất mát và quy trình cập nhật.} Với logit $z_{ik}$:
\begin{equation}
 p_{ik}=\frac{\exp(z_{ik})}{\sum_{j=1}^{10}\exp(z_{ij})},\qquad
 \mathcal{L}_{CE}=-\frac{1}{B}\sum_{i=1}^{B}\log p_{i,y_i}.
\end{equation}
Mỗi mini-batch được chuẩn hóa, đưa qua mạng, tính cross-entropy, lan truyền
ngược và cập nhật bằng AdamW. Đạo hàm loss theo logit của một mẫu tỉ lệ với
$p_{ik}-\mathbf{1}[k=y_i]$, nên lớp sai có điểm cao nhận tín hiệu điều chỉnh
mạnh. Gradient norm được giới hạn ở 1,0 trước bước optimizer; weight decay
$10^{-4}$ giúp kiểm soát độ lớn trọng số. Các cơ chế này không bảo đảm hết
quá khớp, vì còn phụ thuộc dữ liệu, kiến trúc và thời gian huấn luyện.

\textbf{Đánh giá.} Sau mỗi epoch, mạng chuyển sang chế độ evaluation để đo
accuracy validation và lưu checkpoint tốt nhất. Khi suy luận, dropout được
vô hiệu hóa, BatchNorm dùng thống kê đã tích lũy; trọng số không cập nhật.
Nhãn cuối là $\hat y=\arg\max_k z_k$. Softmax chỉ chuyển logit thành phân bố
điểm tổng bằng 1; điểm cao không bảo đảm dự đoán đúng hay xác suất đã được hiệu chuẩn.

\subsection{AlexNet với GroupNorm}
\textbf{Ý tưởng kiến trúc.} AlexNet sử dụng tích chập, ReLU, pooling và
bộ phân lớp có dropout \cite{krizhevsky2012alexnet}. Biến thể CIFAR-10 gồm
năm convolution, với GroupNorm sau mỗi lớp. Các convolution có stride 1;
padding 2 cho kernel $5\times5$ và padding 1 cho kernel $3\times3$.
Vì vậy, chỉ các lớp pooling làm giảm kích thước không gian.

\begin{table}[H]
\centering\small
\caption{Luồng tensor của AlexNet trong thí nghiệm. Conv được theo sau bởi GroupNorm và ReLU.}
\begin{tabular}{lll}
\toprule
Bước & Thao tác & Đầu ra \\
\midrule
Đầu vào & RGB đã chuẩn hóa & $3\times32\times32$ \\
1 & Conv $5\times5$, 64 kênh; max-pool & $64\times16\times16$ \\
2 & Conv $3\times3$, 192 kênh; max-pool & $192\times8\times8$ \\
3 & Conv $3\times3$, 384 kênh & $384\times8\times8$ \\
4 & Conv $3\times3$, 256 kênh & $256\times8\times8$ \\
5 & Conv $3\times3$, 256 kênh; max-pool & $256\times4\times4$ \\
6 & Adaptive average pooling $2\times2$; flatten & $1024$ \\
7 & Dropout; Linear $1024\rightarrow512$; ReLU & $512$ \\
8 & Dropout; Linear $512\rightarrow10$ & $10$ logit \\
\bottomrule
\end{tabular}
\end{table}

\textbf{GroupNorm.} Tại mỗi mẫu, các kênh được chia thành tám nhóm. Trung bình
và phương sai được tính trên các kênh và vị trí không gian trong từng nhóm,
không trên batch \cite{wu2018groupnorm}. Nếu $\mathcal{G}$ là nhóm chứa $a_j$:
\begin{equation}
 \hat a_j=\gamma_j\frac{a_j-\mu_{\mathcal G}}
 {\sqrt{\sigma_{\mathcal G}^2+\epsilon}}+\beta_j.
\end{equation}
Với lớp 64 kênh, mỗi nhóm có tám kênh; với lớp 192 kênh, mỗi nhóm có 24 kênh.
Hệ số scale và bias được học cùng mạng. GroupNorm không cần running mean
và running variance như BatchNorm, nên cách tính thống kê không phụ thuộc số
mẫu trong mini-batch. Không thể từ đó kết luận GroupNorm là nguyên nhân duy nhất
khiến AlexNet đạt kết quả cao nhất; cần thí nghiệm thay riêng lớp chuẩn hóa.

\textbf{Bộ phân lớp và điều chuẩn.} Adaptive pooling giữ lưới $2\times2$,
cho phép vector cuối còn thông tin vị trí thô. Hai lớp tuyến tính học tổ hợp
phi tuyến của các đặc trưng tích chập. Dropout 0,5 ngẫu nhiên loại bỏ một nửa
phần tử ở vị trí dropout khi train và điều chỉnh tỉ lệ phần tử còn lại; khi
inference dùng toàn bộ phần tử. Cơ chế này hạn chế việc bộ phân lớp phụ thuộc
quá mạnh vào một số đặc trưng. Mạng có 2.786.890 tham số. So với AlexNet
ImageNet, ảnh đầu vào, kernel, pooling, chuẩn hóa và classifier đều đã được
điều chỉnh; vì vậy tên kiến trúc cần được đọc cùng cấu hình trên.

\subsection{VGG11-BN với số kênh giảm}
\textbf{Ý tưởng kiến trúc.} VGG dùng các convolution $3\times3$ xếp chồng
\cite{simonyan2015vgg}. Bản cài đặt dùng tám convolution, mỗi lớp có padding 1,
stride 1, theo sau bởi BatchNorm và ReLU. Hai lớp $3\times3$ liên tiếp có vùng
nhìn danh nghĩa $5\times5$ trong trường hợp stride 1, đồng thời bổ sung một
phép phi tuyến giữa chúng. Năm pooling giảm dần kích thước ảnh về $1\times1$.

\begin{table}[H]
\centering\small
\caption{Năm stage của VGG11-BN. Mỗi stage kết thúc bằng max-pooling $2\times2$.}
\begin{tabular}{lllr}
\toprule
Stage & Các convolution $3\times3$ & Đầu ra sau pooling & Số conv \\
\midrule
1 & $3\rightarrow32$ & $32\times16\times16$ & 1 \\
2 & $32\rightarrow64$ & $64\times8\times8$ & 1 \\
3 & $64\rightarrow128\rightarrow128$ & $128\times4\times4$ & 2 \\
4 & $128\rightarrow256\rightarrow256$ & $256\times2\times2$ & 2 \\
5 & $256\rightarrow256\rightarrow256$ & $256\times1\times1$ & 2 \\
\bottomrule
\end{tabular}
\end{table}

\textbf{BatchNorm.} Với mỗi kênh, BatchNorm chuẩn hóa activation bằng trung
bình và phương sai của các mẫu và vị trí không gian trong mini-batch lúc train,
rồi học scale và bias \cite{ioffe2015batchnorm}. Khi evaluation, lớp này dùng
thống kê tích lũy thay vì tính lại theo batch test. Điều này làm thao tác train
và inference khác nhau; chế độ \texttt{model.eval()} là cần thiết để suy luận
đúng. Khác với GroupNorm của AlexNet, thống kê lúc train phụ thuộc batch.

\textbf{Phân lớp và đánh đổi.} Sau stage 5, flatten tạo 256 đặc trưng, đưa
thẳng vào lớp tuyến tính $256\rightarrow10$; không dùng các fully connected lớn
của VGG nguyên bản và không có dropout trong cài đặt này. Mạng có 2.310.186
tham số, ít nhất trong bốn mạng. Cấu trúc tuần tự dễ triển khai nhưng việc giảm
$32\rightarrow16\rightarrow8\rightarrow4\rightarrow2\rightarrow1$ có thể
làm mất chi tiết ở ảnh nhỏ. Dữ liệu không xác nhận riêng tác động của pooling;
đây là một yếu tố cần kiểm tra khi giải thích chênh lệch với các mạng khác.

\subsection{ResNet18 với kết nối phần dư}
\textbf{BasicBlock.} Mỗi block chứa hai convolution $3\times3$ kèm BatchNorm;
ReLU đặt sau convolution thứ nhất và sau phép cộng nhánh tắt. Viết chi tiết:
\begin{align}
 F(x)&=\mathrm{BN}_2\bigl(\mathrm{Conv}_2(
       \mathrm{ReLU}(\mathrm{BN}_1(\mathrm{Conv}_1(x))))\bigr),\\
 y&=\mathrm{ReLU}\bigl(F(x)+S(x)\bigr).
\end{align}
Khi số kênh và kích thước không đổi, $S(x)=x$. Khi đổi kênh hoặc stride 2,
$S$ là convolution $1\times1$ với cùng stride và BatchNorm, để hai nhánh có
cùng shape trước phép cộng. Hai convolution trong nhánh $F$ không có bias
vì đã đi kèm BatchNorm.

\textbf{Vì sao học phần dư?} Thay vì học toàn bộ ánh xạ $H(x)$, nhánh biến đổi
học $F(x)=H(x)-x$ khi dùng identity shortcut \cite{he2016resnet}.
Bỏ qua ReLU để xét phép cộng, đạo hàm theo đầu vào là
$\partial(F(x)+x)/\partial x=\partial F/\partial x+I$.
Hạng $I$ tạo đường truyền trực tiếp cho gradient. ReLU sau phép cộng vẫn có
thể chặn một số thành phần; shortcut giúp tối ưu thuận lợi hơn nhưng không
bảo đảm mọi gradient luôn được giữ nguyên hoặc mạng sâu luôn tốt hơn.

\begin{table}[H]
\centering\small
\caption{Các stage của ResNet18 CIFAR; mỗi stage có hai BasicBlock.}
\begin{tabular}{llll}
\toprule
Phần & Số kênh & Stride block đầu & Đầu ra \\
\midrule
Stem: Conv $3\times3$ + BN + ReLU & 32 & 1 & $32\times32\times32$ \\
Stage 1 & 32 & 1 & $32\times32\times32$ \\
Stage 2 & 64 & 2 & $64\times16\times16$ \\
Stage 3 & 128 & 2 & $128\times8\times8$ \\
Stage 4 & 256 & 2 & $256\times4\times4$ \\
Global average pooling; linear & 256 & -- & $256\rightarrow10$ \\
\bottomrule
\end{tabular}
\end{table}

\textbf{Điều chỉnh cho ảnh nhỏ.} Stem dùng kernel $3\times3$, stride 1,
không max-pooling đầu mạng để giữ chi tiết ở độ phân giải $32\times32$.
Tám BasicBlock có tổng 16 convolution trên nhánh chính; thêm convolution stem
và lớp tuyến tính cuối tạo cấu trúc ResNet18 thông dụng (không đếm riêng
convolution projection trong cách gọi này). Base width giảm còn 32, nên mạng
có 2.797.610 tham số. Global average pooling lấy trung bình trên $4\times4$
cho mỗi kênh, giảm số tham số classifier. ResNet có đường truyền gradient
thuận lợi, nhưng mức accuracy cuối vẫn phụ thuộc cách chuẩn hóa, điều chuẩn,
lịch learning rate và khả năng tổng quát hóa trên dữ liệu đã chọn.

\subsection{Vision Transformer: ViT tiny}
\textbf{Bước 1: biến ảnh thành token.} Cài đặt dùng convolution kernel 4,
stride 4 để nhúng patch. Ảnh $3\times32\times32$ tạo lưới $8\times8$ patch;
mỗi patch có $4\times4\times3=48$ giá trị và được ánh xạ thành vector 192 chiều.
Flatten lưới và chuyển trục tạo tensor $64\times192$. Thêm class token và
embedding vị trí học được:
\begin{equation}
 Z_0=[x_{cls};\,x_p^1E;\ldots;x_p^{64}E]+E_{pos},\quad
 E\in\mathbb{R}^{48\times192},\quad E_{pos}\in\mathbb{R}^{65\times192}.
\end{equation}
Phép chiếu patch thực tế còn có bias. Class token là một vector học được,
không phải ảnh hay nhãn thật; nó thu thập thông tin cho quyết định phân lớp.
Embedding vị trí phân biệt các vị trí trong lưới, vì attention đơn thuần không
mã hóa thứ tự không gian \cite{dosovitskiy2021vit}.

\textbf{Bước 2: self-attention ba head.} Trong mỗi block, LayerNorm chuẩn hóa
192 chiều của từng token; một lớp tuyến tính tạo $Q,K,V$. Mỗi head có
$Q_h,K_h,V_h\in\mathbb{R}^{65\times64}$:
\begin{equation}
 A_h=\operatorname{softmax}_{\mathrm{hàng}}\left(
 \frac{Q_hK_h^\top}{\sqrt{64}}\right),\qquad U_h=A_hV_h.
\end{equation}
Mỗi hàng của $A_h\in\mathbb{R}^{65\times65}$ biểu diễn trọng số kết hợp các
token cho một token truy vấn. Ba $U_h$ được nối thành 192 chiều rồi chiếu
bằng lớp tuyến tính. Các head có thể học những kiểu quan hệ khác nhau, nhưng
không có ràng buộc rằng mỗi head phải tương ứng một bộ phận vật thể cụ thể.
Việc trao đổi thông tin diễn ra giữa mọi cặp token ngay trong một block,
khác với tích chập chỉ kết hợp lân cận tại từng lớp.

\textbf{Bước 3: MLP và kết nối phần dư.} Sáu block dùng cấu trúc pre-norm,
đúng theo thứ tự trong mã nguồn:
\begin{align}
 Z'_\ell&=Z_{\ell-1}+\mathrm{MSA}(\mathrm{LN}(Z_{\ell-1})),\\
 Z_\ell&=Z'_\ell+\mathrm{MLP}(\mathrm{LN}(Z'_\ell)),\\
 \mathrm{MLP}(u)&=W_2\,\mathrm{GELU}(W_1u+b_1)+b_2.
\end{align}
MLP mở rộng $192\rightarrow768$ rồi thu về $192$, xử lý từng token với cùng
trọng số; self-attention đảm nhiệm trao đổi giữa token. LayerNorm không dùng
thống kê batch hay running statistics. Shape $65\times192$ không đổi qua
sáu block; hai nhánh cộng phần dư giúp giữ thông tin đầu vào block.

\begin{table}[H]
\centering\small
\caption{Luồng tensor của ViT tiny trong cài đặt.}
\begin{tabular}{lll}
\toprule
Bước & Cấu hình & Đầu ra \\
\midrule
Patch embedding & Conv kernel 4, stride 4, 192 kênh & $192\times8\times8$ \\
Chuỗi token & Flatten lưới; thêm class token và vị trí & $65\times192$ \\
Encoder & 6 block; 3 head; MLP ratio 4 & $65\times192$ \\
Đầu phân lớp & LayerNorm; lấy class token; Linear & $192\rightarrow10$ \\
\bottomrule
\end{tabular}
\end{table}

\textbf{Bước 4: phân lớp và giới hạn.} Sau encoder, LayerNorm được áp dụng
cho biểu diễn class token trước lớp tuyến tính 10 logit. Mạng có 2.693.578
tham số, không có dropout trong các block và không dùng pretrained.
Chi phí tạo attention tăng theo $T^2$ với $T=65$; mỗi head có 4.225 trọng số
attention cho một ảnh. Patch lớn làm chuỗi ngắn hơn nhưng gom nhiều pixel vào
một token; patch nhỏ tăng số token và chi phí. Trong cấu hình này, ViT có
khả năng mô hình hóa quan hệ toàn ảnh nhưng phải học từ lượng dữ liệu tương
đối hạn chế. Cấu trúc đó không tự bảo đảm accuracy cao hơn CNN.

\subsection{Các phương pháp truyền thống làm mốc so sánh}
Pixel được trải phẳng thành 3.072 chiều, giữ thứ tự các vị trí RGB nhưng không
có cơ chế tích chập hay chia sẻ trọng số theo không gian. HOG chia ảnh xám thành
cell $4\times4$, tính histogram 9 hướng gradient, chuẩn hóa trên block
$2\times2$ chồng lấn và nối thành 1.764 chiều \cite{dalal2005hog}.
HOG giữ bố cục gradient ở lưới cố định, nhưng không giữ trực tiếp màu RGB.

SIFT phẳng tìm keypoint trong một octave, gán hướng và mô tả mỗi điểm bằng
128 giá trị; tối đa 64 descriptor được xếp theo độ mạnh, nối và đệm số 0,
thành 8.192 chiều \cite{lowe2004sift}. Cài đặt này rút gọn SIFT và không dùng
Bag of Visual Words. Khi keypoint thay đổi thứ tự giữa các ảnh, một chiều
trong vector không còn chắc chắn mang cùng ý nghĩa; việc phóng ảnh nhỏ lên
cũng không tạo thêm chi tiết thị giác thực sự.

SVM tuyến tính học các điểm số $s_k=w_k^\top x+b_k$ bằng hinge loss đa lớp
và điều chuẩn L2 \cite{cortes1995svm}. Dù có nhiều chiều, ranh giới trong
không gian đặc trưng vẫn tuyến tính. Random Forest dùng nhiều cây học từ mẫu
bootstrap, chọn ngẫu nhiên đặc trưng khi chia nút và bỏ phiếu nhãn
\cite{breiman2001random}. Các cây tạo ranh giới phi tuyến nhưng không có
cơ chế chuyên biệt để ghép mẫu cục bộ thành biểu diễn ảnh nhiều tầng.
Sáu kết quả của hai bộ phân lớp trên ba loại đầu vào được giữ nguyên trong
bảng so sánh; chúng là baseline cho cấu hình học sâu đã khảo sát.
